\scanpage{24}
% loc 1, heading
$A$ algèbre d'un $k$-\deleted{groupe}\inserted{monoïde} affine $G$

% loc 2
$M$ $k$ module

% loc 3–4, opération and brace
opération de $G$ on $M$ :
\[
 \text{i.e.}\quad G\longrightarrow
 \underline{\unsure{\mathrm{End}|\mathbb E\mathrm{nd}}}_{k}(M)
 =\underline{\mathbb E\mathrm{nd}}_{\unsure{\mathcal O|O|\mathbb O}}\overset{\deleted{\unsure{gxx|gx}}}{(W(M))}
 \quad
 \left\{\begin{aligned}1)&\ 1x=x\\2)&\ g(g'x)=gg'x\end{aligned}\right.
\]

% loc 5–8, Opération sans axiomes
``Opération sans axiomes'' correspond à
\[
\begin{aligned}
 &\deleted{\unsure{u:A\otimes M_{A}\to A\otimes_{k}M}}\quad
 u:M_A\mathrel{\mathop{\unsure{\text{---}|\to}}\limits_{\text{$A$ lin}}}M_A
 \qquad M_A\underset{\unsure{,|1}}{=}\tikz[baseline=(t.base)]{\node[inner sep=0pt](t){$A\otimes_k M$};\draw[thin](-.08,-.31) .. controls (.12,-.38) and (.05,.09) .. (.22,.21) .. controls (.38,.38) and (.58,.35) .. (.64,.30);}\\
 &\deleted{Ax} \\
 &\hspace{4em}\text{ou}\qquad u_0:M\longrightarrow
 \tikz[baseline=(t.base)]{\node[inner sep=0pt](t){$A\otimes_k M$};\draw[thin](-.12,-.30) .. controls (.10,-.39) and (.06,.06) .. (.25,.23) .. controls (.44,.35) and (.60,.31) .. (.64,.24);}
 \qquad\text{$k$-lin.}
\end{aligned}
\]

% loc 9, cancelled sentence
\deleted{L'axiome $1x=x$}

% loc 10–13, Si g
Si $g\in G(k')$ correspond à un hom de $k$-algèbres $A\overset{g}{\longrightarrow}k'$, alors $g_M$ sera l'endomorphisme \deleted{$\unsure{k'\otimes_A}$}\inserted{de} \tikz[baseline=(t.base)]{\node[inner sep=0pt](t){$(\deleted{g}k',\deleted{k}g)\otimes_A M$};\draw[thin](.18,-.29) .. controls (.35,-.33) and (.37,.08) .. (.50,.23) .. controls (.68,.43) and (.81,.36) .. (.81,.31);} \deleted{défini}\inserted{déduit} \deleted{\unsure{comme}}\inserted{de $u$} \deleted{$g\otimes_A i$} par changement de base, i.e. l'unique hom $M_{k'}\longrightarrow M_{k'}$ tel que le diagramme d'hom $A$-linéaires
% loc 14, bracketed fork and square
\diagram{\begin{tikzcd}[ampersand replacement=\&,column sep=2.9em,row sep=1.1em]
 \& M_A\arrow[rr,"u"]\arrow[dd,"{\tikz[baseline=(t.base)]{\node[inner sep=0pt](t){$\scriptstyle g\otimes\mathrm{id}_M$};\draw[thin](-.13,-.16) .. controls (-.06,-.19) and (-.09,.07) .. (.02,.16) .. controls (.13,.30) and (.22,.25) .. (.23,.22);}}"] \& \& M_A\arrow[dd,"{\tikz[baseline=(t.base)]{\node[inner sep=0pt](t){$\scriptstyle g\otimes\mathrm{id}_M$};\draw[thin](-.13,-.16) .. controls (-.06,-.19) and (-.09,.07) .. (.02,.16) .. controls (.13,.30) and (.22,.25) .. (.23,.22);}}"]\\
 \left[\vphantom{\dfrac{\dfrac{A}{A}}{\dfrac A A}}\right. M
 \arrow[ur]\arrow[dr] \&[-1.7em] \hspace{-2.2em}\left.\vphantom{\dfrac{\dfrac{A}{A}}{\dfrac A A}}\right] \& \& \\
 \& M_{k'}\arrow[rr,swap,"g_M"] \& \& M_{k'}
\end{tikzcd}}
% loc 15, after first diagram
est commutatif ; Mais cela signifie aussi que le diagramme
% loc 16, square
\diagram{\begin{tikzcd}[ampersand replacement=\&]
 M\arrow[r,"u_0"]\arrow[d,swap,"\mathrm{can}"] \& M\deleted{\unsure{{}_A}}\otimes A\arrow[d,"\mathrm{id}_M\otimes g"]\\
 M\deleted{\unsure{{}_k|{}_{k'}}}\otimes k'\arrow[r,swap,"g_M"] \& M\otimes k'
\end{tikzcd}}
% loc 17–19, est commutatif
est commutatif, ou encore, notant que la donnée de $g_M$ est équivalente à celle de $g_M^{\unsure{0|\circ}}=g_M\circ\mathrm{can}:M\mathrel{\tikz[baseline=-.3ex,x=1em,y=1em]{\draw[line width=.45pt] (0,0)--(1,0)--(1,.16);}}M\otimes k'$, à la \deleted{\unsure{fon.}} commutativité du
% loc 20, triangle
\diagram{\begin{tikzcd}[ampersand replacement=\&]
 M\arrow[r,"u_0"]\arrow[dr,swap,"g_M^{\unsure{0|\circ}}"] \& M\otimes A\arrow[d,"\unsure{\mathrm{id}|\mathrm{Id}}_M\otimes g"]\\
 \& M\otimes k'
\end{tikzcd}}

% loc 21–22, Pour traduire
Pour traduire l'axiome 1), on fait $k'=k$, $g=\deleted{\mathrm{id}_k}$ $(\varepsilon:A\longrightarrow k)$, et on trouve
% loc 23, circled 1 and triangle with equation
\diagram{\begin{tikzcd}[ampersand replacement=\&,column sep=2.8em,row sep=large]
 \tikz[baseline=(t.base)]{\node[draw,circle,inner sep=1.5pt](t){$1'$};} \& M\arrow[r,"u_0"]\arrow[dr,swap,near start,"\text{---}"{sloped}] \& M\otimes A\arrow[d,"\mathrm{id}_M\otimes\varepsilon"] \& \text{comm. i.e.}\\
 \& \& M \& (\mathrm{id}_M\otimes\varepsilon)u_0=\mathrm{id}_M
\end{tikzcd}}
% archive sheet number: 23
